\chapter{Journalisation, capture et rejeu}\label{ch:ll:logging-capture-and-replay}

Un régulateur demande pourquoi un ordre est parti à 09:30:00.012345678 avec ce prix et cette taille. Il y a deux façons de répondre. L'une
consiste à chercher dans les journaux les lignes proches de cet instant et à reconstituer, à partir de ce que les programmeurs ont choisi d'afficher, ce que la stratégie
a probablement vu. L'autre consiste à prendre les entrées que le processus de trading a reçues ce matin-là, chacune avec l'instant de sa réception,
à les fournir au même binaire sur un poste de travail, et à s'arrêter sur l'ordre~: le carnet, les paramètres en vigueur et chaque valeur utilisée par la
décision sont là, exactement tels qu'ils étaient, parce que le programme est déterministe (chapitre~20) et que ses entrées ont été conservées.
Ce chapitre construit la seconde réponse et la machinerie dont elle a besoin~: un journal qui coûte quelques nanosecondes sur le hot path parce
qu'il écrit des nombres bruts et les formate plus tard, un \omterm{def:ll:logging-capture-and-replay:journal}{journal des entrées} qui enregistre ce qui est arrivé et quand, et un pilote de rejeu qui
fait d'une journée de production une expérience reproductible.

\section{Pourquoi la journalisation formatée n'a pas sa place sur le hot path}

Une ligne de journal formatée effectue un travail surprenant~: elle analyse une chaîne de format, convertit chaque nombre en texte décimal,
recopie le résultat dans un tampon et, selon la bibliothèque, prend un verrou et effectue un appel système. Sur l'ordinateur portable, une ligne
de quatre nombres coûte environ \qty{100}{\nano\second} en médiane avec \texttt{snprintf} vers un tampon, à peu près autant
avec \texttt{fprintf} vers un fichier entièrement bufferisé, et \qty{210}{\nano\second} avec un flux de chaîne
(la \cref{fig:ll:logging-capture-and-replay:calls})~; et sa queue de distribution est longue, parce que le tampon du fichier se remplit et est écrit, et
parce que le flux alloue de la mémoire. Un moteur de stratégie qui traite un événement en \qty{50}{\nano\second} (chapitre~20) ne peut pas se permettre une
telle ligne par événement, et un système dont la journalisation est désactivée en production pour gagner ce temps n'a rien à examiner quand
quelque chose tourne mal.

La réponse n'est pas de journaliser moins, mais de faire moins en journalisant. Tout ce que contient une ligne de journal est connu à la compilation
du programme, sauf les nombres~: le texte, la position des arguments et leurs types. Le hot path peut donc n'écrire que les nombres,
avec un identifiant de l'instruction, et laisser le formatage à un programme qui s'exécute plus tard, ailleurs, quand quelqu'un lit
le journal. La bibliothèque standard \texttt{C++} du compilateur utilisé ici (g++ 11) ne fournit même pas \texttt{std::format}~; cela ne
changerait pas la conclusion, puisque c'est précisément le formatage dont on supprime le coût.

\section{Journalisation binaire et formatage différé}

\begin{definition}[Journal binaire, formatage différé]\label{def:ll:logging-capture-and-replay:binlog}
Un \emph{journal binaire}\index{journal binaire} est un journal dont les enregistrements contiennent un identifiant d'instruction, un horodatage et les valeurs brutes des
arguments de l'instruction, et non du texte~; les chaînes de format des instructions sont stockées une seule fois, dans l'en-tête du journal ou à côté du binaire.
Le \emph{formatage différé}\index{formatage différé} est la pratique qui consiste à ne transformer ces enregistrements en texte qu'au moment où le journal est
lu, par un programme distinct, de sorte que le programme qui écrit le journal ne convertit jamais un nombre en texte.
\end{definition}

L'idée est établie~: NanoLog, présenté à l'USENIX Annual Technical Conference de 2018, extrait à la compilation les parties statiques des instructions
de journalisation, écrit à l'exécution un \omterm{def:ll:logging-capture-and-replay:binlog}{journal binaire} compacté qu'il redéploie hors ligne, et annonce un coût
d'appel de 8 nanosecondes en \omterm{def:ll:cpu-microarchitecture:ubench}{microbenchmarks} et de 18 en application. Le \texttt{firm.binlog} de la construction fait de même avec le
seul langage. Une instruction s'écrit \texttt{log.log<\textquotedbl{}fill \{\} of \{\} at \{\}, position \{\}\textquotedbl{}>(ts, id, qty, price, pos)}~: la chaîne
de format et les types des arguments sont des arguments de template, de sorte que l'identifiant de l'instruction est calculé à la compilation, comme un hachage des deux
(\cref{lst:ll:logging-capture-and-replay:id}), et que l'instruction enregistre son format dans une table pendant l'initialisation statique,
avant que \texttt{main} ne s'exécute. À l'exécution, un appel écrit l'identifiant, l'horodatage et les octets des arguments dans un emplacement d'un
anneau à producteur unique (chapitre~12), puis rend la main (\cref{lst:ll:logging-capture-and-replay:log}). Si l'anneau est plein, l'enregistrement
est abandonné et compté~: le hot path n'attend jamais le journal.

\omcode{code/firm/binlog/cpp/firm_binlog.hpp}{62}{70}{L'identifiant d'une instruction, calculé par le compilateur à partir de son format et
des types de ses arguments.}\label{lst:ll:logging-capture-and-replay:id}

\omcode{code/firm/binlog/cpp/firm_binlog.hpp}{125}{141}{Le hot path d'un appel de journalisation~: encoder les arguments bruts derrière un
en-tête de 16 octets et écrire l'enregistrement dans l'anneau, ou l'abandonner.}\label{lst:ll:logging-capture-and-replay:log}

Un thread d'écriture, sur un autre cœur, vide l'anneau dans un fichier. Le fichier commence par un en-tête qui liste chaque instruction
(identifiant, types des arguments, format), triées par identifiant, puis les enregistrements dans l'ordre où ils sont sortis de l'anneau~; le décodeur Python
lit l'en-tête, décode chaque enregistrement avec les types de son instruction et le formate. Le format est assez simple pour être écrit par
trois implémentations~: l'encodeur de référence Python écrit le journal d'exemple de la fixture, et le logger C++, par son anneau et
par son thread d'écriture, ainsi que le logger Rust écrivent chacun les mêmes instructions avec les mêmes octets. Un test compte aussi les allocations
sur le tas pendant la journalisation et n'en trouve aucune, et remplit volontairement un petit anneau pour vérifier que le débordement est abandonné et compté.

\begin{omfigure}
\begin{tikzpicture}[font=\scriptsize,>=Stealth,box/.style={draw,rounded corners=2pt,minimum height=0.75cm,align=center,
  inner sep=3pt,fill=omDef!12}]
\node[box] (hot) at (0,0) {hot path\\ \texttt{log<\textquotedbl{}fill \{\} of \{\}\textquotedbl{}>}};
\node[box,fill=omMeth!15,minimum width=2.4cm] (ring) at (3.8,0) {anneau (SPSC)\\ id $|$ len $|$ ts $|$\\ arguments bruts};
\node[box] (wr) at (7.1,0) {thread d'écriture\\ (autre cœur)};
\node[box,fill=gray!10] (file) at (10.3,0) {fichier~:\\ en-tête des instructions\\ $+$ enregistrements};
\node[box,fill=omThm!15] (dec) at (10.3,-1.9) {décodeur (Python)~:\\ texte, hors ligne};
\node[box] (feed) at (0,-1.9) {entrées telles\\ que reçues (flux,\\ comptes rendus,\\ paramètres)};
\node[box,fill=gray!10] (jr) at (3.5,-1.9) {journal des entrées\\ recv ts $|$ événement};
\node[box,fill=omThm!15] (rep) at (6.3,-1.9) {rejeu~: même\\ binaire, mêmes\\ entrées};
\draw[->,thick] (hot) -- node[above]{\qty{5}{\nano\second}} (ring);
\draw[->] (ring) -- (wr);
\draw[->] (wr) -- (file);
\draw[->] (file) -- (dec);
\draw[->] (feed) -- (jr);
\draw[->] (jr) -- (rep);
\draw[<->,dashed] (rep) -- node[above=1pt,pos=0.45]{comparer} (dec);
\end{tikzpicture}

\omcaption{Journalisation, capture et rejeu. Le hot path écrit des enregistrements bruts dans un anneau et rend la main~; un thread d'écriture les vide
dans un fichier dont l'en-tête contient le format de chaque instruction~; le formatage a lieu hors ligne. Les entrées sont journalisées à leur arrivée,
et un rejeu du journal par le même binaire produit le même journal.}\label{fig:ll:logging-capture-and-replay:flow}
\end{omfigure}

La \cref{fig:ll:logging-capture-and-replay:calls} en donne le résultat~: environ \qty{5}{\nano\second} par appel en médiane, vingt
fois moins que \texttt{snprintf}, avec une queue au 99\up{e}~centile quand le thread d'écriture vide l'anneau et que les \omterm{def:ll:memory-hierarchy-and-caches:line}{lignes de cache} qu'il
lit passent d'un cœur à l'autre. Les manières du thread d'écriture comptent. La première version interrogeait l'anneau sans pause, et chaque appel de journalisation
payait les \omterm{def:ll:memory-hierarchy-and-caches:line}{lignes de cache} que le thread d'écriture ne cessait de lui reprendre~; celui de la construction vide par lots et dort
\qty{100}{\micro\second} quand il y a peu à vider (\cref{lst:ll:logging-capture-and-replay:writer}). Le prix n'est payé par
personne sur le hot path~: le journal atteint le fichier une fraction de milliseconde plus tard.

\omcode{code/firm/binlog/cpp/firm_binlog.hpp}{206}{212}{La boucle du thread d'écriture~: vider par lots, dormir quand il y a peu à
vider.}\label{lst:ll:logging-capture-and-replay:writer}

\begin{omfigure}
\begin{tikzpicture}
\begin{axis}[width=10.5cm,height=5.6cm,ybar,bar width=12pt,ymode=log,ymin=1,ymax=2000,log origin=infty,
  xmin=-0.6,xmax=3.6,xtick={0,1,2,3},xticklabels={firm.binlog,snprintf,fprintf,ostringstream},
  ylabel={ns par appel},tick label style={font=\small},label style={font=\small},grid=major,grid style={gray!25},
  legend pos=north west,legend style={font=\scriptsize},table/col sep=comma]
\addplot[fill=omThm!45,draw=omThm] table[x expr={floor(\coordindex/3)},y=ns,
  restrict expr to domain={mod(\coordindex,3)}{0:0}]{figdata/low-latency/23-logging-capture-and-replay/measured_logcall.csv};
\addlegendentry{médiane}
\addplot[fill=omMeth!40,draw=omMeth] table[x expr={floor(\coordindex/3)},y=ns,
  restrict expr to domain={mod(\coordindex,3)}{1:1}]{figdata/low-latency/23-logging-capture-and-replay/measured_logcall.csv};
\addlegendentry{99\up{e}~centile}
\end{axis}
\end{tikzpicture}

\omcaption{Le coût d'un appel de journalisation portant quatre nombres, selon la méthode (chronométrage au TSC, 200\,000 appels chacune)~: le \omterm{def:ll:logging-capture-and-replay:binlog}{journal binaire} avec
son thread d'écriture vidant l'anneau sur un autre cœur, \texttt{snprintf} vers un tampon, \texttt{fprintf} vers un
fichier entièrement tamponné, et un flux de chaînes. Mesuré sur un ordinateur portable (Intel Core Ultra 7 155H) sous WSL2. Données~:
\texttt{bench\_log.py}.}\label{fig:ll:logging-capture-and-replay:calls}
\end{omfigure}

\section{Capture~: entrées, sorties et temps}

\begin{definition}[Journal des entrées]\label{def:ll:logging-capture-and-replay:journal}
Un \emph{journal des entrées}\index{journal des entrées} est l'enregistrement, tenu par un processus, de chaque entrée reçue, dans l'ordre
où il les a reçues, chacune avec l'heure de sa réception~: événements de marché, comptes rendus d'ordres, déclenchements de minuteries dépendant de l'horloge,
changements de paramètres et messages de contrôle. Avec la version du programme, il suffit pour reproduire tout ce que le
processus a calculé.
\end{definition}

Un journal des décisions dit ce que le programme a fait~; seul un \omterm{def:ll:logging-capture-and-replay:journal}{journal des entrées} dit pourquoi, parce qu'il permet de refaire la décision.
Le journal de la construction est volontairement austère~: un court en-tête, puis des enregistrements fixes de 56~octets, chacun formé de l'heure de réception suivie de
l'événement de 48~octets du \omterm{def:ll:the-feed-handler:handler}{feed handler} (chapitre~18). Il est écrit par le thread qui reçoit les entrées, au point où elles
entrent dans le moteur, de sorte que son ordre soit celui du moteur. Deux pièges sont faciles à commettre. Le journal doit contenir ce que le
processus a reçu, et non ce que la bourse a envoyé~: l'enregistrement de la bourse ne présente ni trous, ni retransmissions, ni horodatages d'arrivée
locaux, et un rejeu de celui-ci répond à une autre question. Les horodatages doivent être ceux du processus lui-même, pris une seule fois, à la réception,
avec l'horloge du chapitre~5, car ce sont aussi des entrées~: une stratégie qui tient compte de l'âge de ses cotations les lit. L'horodatage de la
bourse et l'horodatage de réception du livre~1 (chapitre~28) sont conservés tous deux, l'un à l'intérieur de l'événement, l'autre devant lui.

\begin{dated}{2026-09}{Durée et précision de conservation des enregistrements}\label{dat:ll:logging-capture-and-replay:records}
Dans l'Union européenne (consulté en septembre 2026), la Commission Delegated Regulation (EU) 2017/589 oblige les sociétés qui se livrent
au trading algorithmique à haute fréquence à enregistrer leurs ordres et à conserver ces enregistrements pendant cinq ans à compter de la soumission de chaque
ordre (article 28). La Commission Delegated Regulation (EU) 2017/574 oblige les membres d'une plateforme de négociation qui recourent à une technique de trading algorithmique à haute fréquence
à conserver leurs horloges internes à moins de 100~microsecondes de l'UTC, avec une granularité des
horodatages d'une microseconde ou mieux, et à réviser leur traçabilité à l'UTC au moins une fois par an.
\end{dated}

Les journaux se compressent bien, car des événements consécutifs partagent l'essentiel de leurs octets~: le même instrument, des prix et des horodatages voisins,
des numéros de séquence qui progressent de un en un. Le \Cref{tab:ll:logging-capture-and-replay:journal} mesure cette compression sur les deux flux de
la construction~: environ 3.8 pour 1 avec le compresseur généraliste zlib à son niveau par défaut, soit environ 15~octets par événement.

\begin{table}[ht]
\centering\small
\begin{tabular}{lrrrrr}
\toprule
flux & événements & durée & journal & compressé & taux \\
\midrule
grand tick (l'ouverture du simulateur) & 189\,875 & \qty{120}{\second} & 10.6~Mo & 2.8~Mo & 3.76 \\
petit tick (synthétique) & 6\,000 & \qty{6}{\milli\second} & 0.34~Mo & 0.08~Mo & 3.96 \\
\bottomrule
\end{tabular}
\caption{Le \omterm{def:ll:logging-capture-and-replay:journal}{journal des entrées} de deux flux d'événements~: 56~octets par événement brut, environ 15 après compression avec zlib au niveau 6. Données~:
\texttt{fig\_journal.py} (déterministe).}\label{tab:ll:logging-capture-and-replay:journal}
\end{table}

\section{Rejeu déterministe d'une journée de production}

\begin{definition}[Rejeu déterministe]\label{def:ll:logging-capture-and-replay:replay}
Un \emph{rejeu déterministe}\index{rejeu déterministe} exécute un programme sur un \omterm{def:ll:logging-capture-and-replay:journal}{journal des entrées}, dans l'ordre du journal, avec le temps
pris dans le journal, et produit bit à bit les sorties que le programme a produites lorsque le journal a été enregistré~; il
exige que le programme soit déterministe (chapitre~20) et que le journal soit complet.
\end{definition}

Le pilote de rejeu lit le journal, reconstruit les événements et les remet au moteur de stratégie avec l'horloge simulée~: le test de la
construction consigne dans un journal la journée à petit tick, la rejoue et obtient l'empreinte de sortie enregistrée au chapitre~20. Un rejeu est aussi
rapide, car rien n'attend le marché~: sur l'ordinateur portable, le moteur rejoue les deux minutes les plus actives du simulateur autour de l'ouverture,
189\,875 événements, lecture du journal comprise, en moins de \qty{50}{\milli\second}, environ 2\,800 fois plus vite qu'ils ne se sont produits. Une journée
se rejoue en quelques secondes, ce qui change ce qu'on peut en faire~: toute modification du code peut être rejouée sur la semaine précédente, et
chaque incident peut être rejoué autant de fois que nécessaire, avec un débogueur attaché.

Lorsque deux exécutions qui devraient concorder ne concordent pas, les journaux disent où. Chaque exécution journalise ses actions~; les deux journaux sont comparés enregistrement par
enregistrement (le \cref{lst:ll:logging-capture-and-replay:diff}), et le premier enregistrement qui diffère désigne le premier événement dont le traitement
a différé. Le test de la construction introduit une différence, un changement de paramètre à l'événement 3\,000 que la seconde exécution reçoit et la première non,
comme le ferait une valeur lue hors du journal~; la comparaison trouve la première action qui diffère,
numéro 540 sur 1\,084, juste après l'événement introduit, et rien avant lui. C'est la procédure du problème du week-end du chapitre~20,
mécanisée~: le journal indique où regarder, et le rejeu permet de regarder autant de fois que nécessaire.

\omcode{code/firm/binlog/firm_binlog.py}{99}{107}{Hors ligne~: mise en forme des enregistrements, et recherche du premier enregistrement où deux journaux
diffèrent.}\label{lst:ll:logging-capture-and-replay:diff}

\section{Ce qu'il faut conserver, et pendant combien de temps}

Trois choses sont conservées. Le \omterm{def:ll:logging-capture-and-replay:journal}{journal des entrées}, car c'est le seul enregistrement à partir duquel tout le reste peut être recalculé~; le
journal des décisions, car il répond à la plupart des questions sans rejeu et vérifie le rejeu lorsqu'on en fait un~; et le programme~:
le binaire exact, sa configuration et les instantanés des paramètres, car un journal rejoué avec une autre construction ne répond
à rien. Les horloges et les durées de conservation du régulateur (l'encadré daté) fixent le minimum~: des horodatages traçables à l'UTC avec
la granularité requise, et cinq ans de registres d'ordres. Le stockage est modeste au regard des données qu'une firme quantitative
conserve de toute façon, comme le calcule le problème du week-end, et le coût de ne pas le conserver est la première phrase de ce chapitre.

\section{Tutoriel~: journal, journal des entrées, rejeu, bissection}\label{tut:ll:logging-capture-and-replay:tut}

\begin{tutorial}
\textbf{Objectif.} Journaliser depuis le chemin critique (hot path) en quelques nanosecondes par appel, consigner les entrées d'une journée, les rejouer bit à bit, et trouver une
différence introduite entre deux exécutions. \textbf{État final~:} la \cref{fig:ll:logging-capture-and-replay:calls}, le tableau du journal,
la vitesse du rejeu, et des tests au vert dans trois langages.
\begin{tutsteps}
\item \textbf{Journalisation.} Le test C++ journalise les énoncés du jeu d'essai par l'anneau et par le thread d'écriture, et compare les
octets avec \texttt{data/sample.blog}, écrit par la référence Python~; le test Rust écrit les mêmes octets~; les fonctions \texttt{read} et \texttt{text} du décodeur Python les transforment en
\texttt{data/sample.txt}.
\item \textbf{Mesure.} \texttt{python bench\_log.py} chronomètre un appel de journalisation selon chaque méthode, ainsi que le rejeu du journal de l'ouverture.
\item \textbf{Journal et rejeu.} Le test consigne dans un journal la journée à petit tick, la rejoue au moyen de \texttt{firm.stratengine}, et
vérifie l'empreinte enregistrée.
\item \textbf{Bissection.} \texttt{ll\_bisect\_test.cpp} exécute deux fois le moteur, introduit une modification dans la seconde exécution, journalise les deux, et
trouve le premier enregistrement qui diffère.
\end{tutsteps}
\textbf{Pour aller plus loin.} Remplacer le changement de paramètre introduit par une lecture de l'horloge murale (le défaut du chapitre~20), et observer la
première différence se déplacer d'une exécution à l'autre~; compresser le journal au fur et à mesure de son écriture, dans le thread d'écriture.
\end{tutorial}

\section{Construction~: le journal binaire et le journal des entrées}\label{bld:ll:logging-capture-and-replay:binlog}

\begin{build}
\bfield{Rôle} Chaque composant du chemin de trading journalise par lui~: les trous du \omterm{def:ll:the-feed-handler:handler}{feed handler} (chapitre~18), les actions du
moteur (chapitre~20), les messages de la passerelle (chapitre~21), les décisions de la \omterm{def:ll:pre-trade-risk-and-kill-switches:gate}{barrière de risque} (chapitre~22)~; son journal est ce que
l'architecture résiliente du chapitre~24 réplique et ce que les tests du chapitre~25 rejouent.
\bfield{Interface} C++20 \texttt{firm::binlog}~: \texttt{Logger(capacity)} avec \texttt{log<\textquotedbl{}format\textquotedbl{}>(ts, arguments)},
\texttt{drain}, \texttt{dropped}~; \texttt{Writer(logger, path)}~; \texttt{header()}~; \texttt{Journal} avec \texttt{add} et
\texttt{records}. Python \texttt{firm\_binlog}~: \texttt{site\_id}, \texttt{Writer}, \texttt{read}, \texttt{text},
\texttt{first\_difference}, \texttt{journal\_write}, \texttt{journal\_read}. Rust \texttt{firm\_binlog}~: \texttt{Logger},
\texttt{read}, \texttt{text}, \texttt{site\_id}.
\bfield{Règles} Aucune mise en forme, aucun verrou, aucun appel système ni allocation sur le chemin critique (hot path)~; un anneau plein rejette et compte~; les énoncés
identifiants issus du format et des types~; le journal est écrit là où les entrées entrent dans le moteur, avec les horodatages de réception~; les
mêmes octets issus des trois implémentations.
\bfield{Tests d'acceptation} \texttt{code/firm/binlog/}~: le jeu d'essai écrit octet pour octet par Python, C++ (anneau et thread d'écriture)
et Rust, et décodé jusqu'à son texte~; énoncés aléatoires soumis à un aller-retour~; identifiants qui dépendent des types, et collisions
refusées~; zéro allocation pendant la journalisation, et débordement compté~; un journal rejoué par le moteur de stratégie jusqu'à l'empreinte
enregistrée~; la première différence entre deux journaux est trouvée.
\bfield{Pour aller plus loin} Compression dans le thread d'écriture~; un énoncé par composant avec une sévérité filtrée à la compilation~;
horodatages issus du TSC avec des enregistrements de calibration dans le journal, convertis hors ligne~; un journal des comptes rendus de la passerelle à côté de
les données de marché.
\end{build}

\begin{omsources}
\item S.~Yang, S.~J.~Park et J.~Ousterhout, ``NanoLog: a nanosecond scale logging system'', \emph{USENIX Annual Technical
Conference}, 2018.
\item Commission Delegated Regulation (EU) 2017/589 (RTS 6), article 28.
\item Commission Delegated Regulation (EU) 2017/574 (RTS 25), article 4 et annexe.
\end{omsources}

\section{Exercices}

\begin{exercise}[$\star$]\label{exo:ll:logging-capture-and-replay:1}
Combien d'octets occupe un enregistrement de l'énoncé \texttt{\textquotedbl{}fill \{\} of \{\} at \{\}, position \{\}\textquotedbl{}} avec quatre arguments de 64~bits,
et combien occupe sa ligne formatée pour les valeurs 40, 100, 999\,900 et $-200$ avec un horodatage de 14~chiffres~? Quelle quantité
le \omterm{def:ll:logging-capture-and-replay:binlog}{journal binaire} économise-t-il alors~?
\end{exercise}

\begin{exercise}[$\star$]\label{exo:ll:logging-capture-and-replay:2}
Pourquoi l'identifiant de l'énoncé dépend-il des types de ses arguments et pas seulement de son format~?
\end{exercise}

\begin{exercise}[$\star$]\label{exo:ll:logging-capture-and-replay:3}
Le journal de l'ouverture compte 189\,875 enregistrements de 56~octets plus un en-tête de 12~octets. Vérifiez sa taille, et les octets par événement
après compression, à partir du \cref{tab:ll:logging-capture-and-replay:journal}.
\end{exercise}

\begin{exercise}[$\star\star$]\label{exo:ll:logging-capture-and-replay:4}
Un composant journalise 2 millions d'enregistrements par seconde pendant une rafale~; le thread d'écriture est bloqué pendant \qty{20}{\milli\second} par
un disque lent. Combien d'emplacements d'anneau faut-il pour qu'aucun enregistrement ne soit perdu, et quelle mémoire cela représente-t-il avec des emplacements de 64~octets~?
\end{exercise}

\begin{exercise}[$\star\star$]\label{exo:ll:logging-capture-and-replay:5}
Le chemin critique (hot path) horodate les enregistrements avec le TSC plutôt qu'avec \texttt{clock\_gettime}. Que doit contenir en plus le journal pour que les
horodatages puissent être convertis en UTC hors ligne avec une granularité de la microseconde~?
\end{exercise}

\begin{exercise}[$\star\star$]\label{exo:ll:logging-capture-and-replay:6}
Un moteur traite un événement en \qty{50}{\nano\second} et journalise cinq lignes par événement. Que coûte la journalisation par événement avec le
\omterm{def:ll:logging-capture-and-replay:binlog}{journal binaire} et avec \texttt{snprintf}, aux médianes de la \cref{fig:ll:logging-capture-and-replay:calls}~?
\end{exercise}

\begin{exercise}[$\star\star\star$]\label{exo:ll:logging-capture-and-replay:7}
\emph{Programmation.} Journalisez les décisions de la \omterm{def:ll:pre-trade-risk-and-kill-switches:gate}{barrière de risque} (risk gate, chapitre~22) avec un énoncé portant l'identifiant de l'ordre, le code et le
masque~; rejouez le jeu d'essai de la barrière, et décodez le journal en un compte des refus par motif qui correspond à \texttt{expected.txt}.
\end{exercise}

\begin{exercise}[$\star\star\star$]\label{exo:ll:logging-capture-and-replay:8}
\emph{Trouver la faille.} «~Nous n'avons pas besoin de \omterm{def:ll:logging-capture-and-replay:journal}{journal des entrées}. Nous journalisons chaque décision avec son motif, et pour rejouer une journée nous faisons tourner la
stratégie sur l'enregistrement propre de la bourse des données de marché.~»
\end{exercise}

\section{Problème~: jusqu'à la nanoseconde}

\begin{problem}[{Problème du week-end --- ce qu'il en coûte pour pouvoir répondre}]\label{pb:ll:logging-capture-and-replay:1}
Une entreprise journalise les entrées de ses processus de trading et doit pouvoir rejouer tout ordre des cinq dernières années. Utilisez le
\cref{tab:ll:logging-capture-and-replay:journal} et la vitesse de rejeu mesurée.

\medskip\noindent\textbf{Partie I --- Le débit.}
\begin{enumerate}
\item Quel est le débit moyen d'événements du flux de l'ouverture~?
\item À ce débit moyen, combien d'événements contient une journée de négociation de 6.5~heures~?
\item Quelle est la taille brute de ce journal en octets, et compressé au taux mesuré~?
\item Pourquoi une moyenne prise sur une ouverture chargée est-elle une estimation pessimiste pour une journée entière, et pourquoi reste-t-elle néanmoins celle qu'il faut retenir
pour planifier~?
\end{enumerate}

\medskip\noindent\textbf{Partie II --- Cinq ans.}
\begin{enumerate}[resume]
\item Avec 252 jours de négociation par an, quel volume occupent cinq ans de tels journaux, compressés~?
\item L'entreprise négocie 200 instruments aux flux similaires. Passer la réponse à l'échelle.
\item Que faut-il encore conserver pour qu'un journal puisse être rejoué cinq ans plus tard~?
\item Pourquoi le journal des décisions, bien que plus petit, ne suffit-il pas~?
\end{enumerate}

\medskip\noindent\textbf{Partie III --- Le rejeu.}
\begin{enumerate}[resume]
\item À la vitesse mesurée, combien de temps dure le rejeu de la journée d'un instrument~?
\item Et celui des 200 instruments, les uns après les autres, sur un seul cœur~?
\item Qu'est-ce qui limite la vitesse de rejeu, et qu'est-ce qui ne la limite pas~?
\item Pourquoi le rejeu doit-il lire le temps dans le journal plutôt que dans la machine~?
\end{enumerate}

\medskip\noindent\textbf{Partie IV --- Le verdict.}
\begin{enumerate}[resume]
\item Énoncer le \emph{résultat nommé}~: le volume de stockage du journal par jour au débit et au taux de compression mesurés, et la vitesse de rejeu
par rapport au temps réel.
\item Que coûte un appel de journalisation avec \omterm{def:ll:logging-capture-and-replay:binlog}{formatage différé}, et avec \texttt{snprintf}~?
\item Que doit reproduire le rejeu pour répondre à la question du régulateur, et comment vérifier qu'il l'a fait~?
\item Que devient la rejouabilité lorsqu'un composant lit l'horloge murale~?
\item À quel endroit du processus le journal doit-il être écrit, et pourquoi là~?
\item Que démontre le test de différence plantée, et que faudrait-il pour trouver une vraie différence~?
\item Quelle précision d'horodatage exige l'encadré daté, et un journal horodaté au TSC la satisfait-il~?
\item En une phrase~: qu'est-ce qui rend un ordre explicable des années plus tard~?
\end{enumerate}
\end{problem}

\section{Questions d'entretien}

\begin{interviewq}[$\star$ \iqroles{développeur}]\label{iq:ll:logging-capture-and-replay:1}
Pourquoi la journalisation de style \texttt{printf} pose-t-elle un problème sur le chemin critique (hot path), et que faites-vous à la place~?
\end{interviewq}

\begin{interviewq}[$\star\star$ \iqroles{développeur}]\label{iq:ll:logging-capture-and-replay:2}
Concevez un journaliseur dont l'appel coûte moins de dix nanosecondes. Que se passe-t-il lorsque son tampon est plein~?
\end{interviewq}

\begin{interviewq}[$\star\star$ \iqroles{développeur}]\label{iq:ll:logging-capture-and-replay:3}
Qu'enregistreriez-vous pour pouvoir rejouer exactement une journée de négociation, et à quel endroit du processus l'enregistreriez-vous~?
\end{interviewq}

\begin{interviewq}[$\star\star$ \iqroles{développeur}]\label{iq:ll:logging-capture-and-replay:4}
Deux rejeux de la même journée ne concordent pas. Comment trouvez-vous où~?
\end{interviewq}

\begin{interviewq}[$\star\star$ \iqroles{développeur, risque}]\label{iq:ll:logging-capture-and-replay:5}
Un régulateur demande pourquoi un ordre a été envoyé il y a trois ans. De quoi avez-vous besoin pour répondre, et combien de temps cela prend-il~?
\end{interviewq}

\begin{interviewq}[$\star\star\star$ \iqroles{développeur}]\label{iq:ll:logging-capture-and-replay:6}
Concevez la capture et le rejeu pour une entreprise qui compte des dizaines de processus de trading sur plusieurs plateformes. Que journalise-t-on, où, dans quel
format, pendant combien de temps, et comment teste-t-on le rejeu~?
\end{interviewq}
